\documentclass[11pt,a4paper]{article}

% ---------- Packages ----------
\usepackage[margin=1in]{geometry}
\usepackage{enumitem}
\usepackage[hidelinks]{hyperref}
\usepackage{titlesec}
\usepackage{tabularx}
\usepackage{setspace}
\usepackage[section]{placeins}
\usepackage{needspace}


% ---------- Formatting ----------
\setstretch{1.05}
\pagenumbering{gobble}
\setlength{\parindent}{0pt}

\titleformat{\section}
  {\large\bfseries}
  {}
  {0em}
  {}[\titlerule]

\titleformat{\subsection}
  {\bfseries}
  {}
  {0em}
  {}

\hypersetup{
  colorlinks=true,
  linkcolor=black,
  urlcolor=blue
}

% ---------- Document ----------
\begin{document}

% ---------- Header ----------
\begin{center}
    {\LARGE \textbf{Tran Gia Khanh}} \\[4pt]
    27 Tran Duy Hung, Hanoi \;|\; Email:\href{mailto:giakhanh28031995@gmail.com}{giakhanh28031995@gmail.com} \;|\;
    0353102005 \\[2pt]
    Linkedin:https://www.linkedin.com/in/khanh-tran-gia-9944a9372/
\end{center}

% ---------- Education ----------
\section*{Education}
\textbf{University of Science and Technology of Hanoi} \hfill 2023 -- 2026 \\
Bachelor of Data Science (Undergraduate) \\
GPA: 3.5 / 4.0

\textbf{Vin University} \hfill April 2026 - Present \\
VinGroup's Applied AI Talent Program(Al Thuc Chien)


% ---------- Experience ----------
\section*{Experience}
\textbf{Ploomba} — Canadian Agricultural Technology Development  Startup \hfill Oct 2025 -- Present \\
Company: {https://www.linkedin.com/company/ploomba/posts/?feedView=all}
\begin{itemize}[leftmargin=1.5em]
    \item Position: Data Assistant, AI engineer.
    \item Built and fine-tuned custom YOLO-based computer vision models for fruit detection in agricultural environments, supporting automated farming systems. Annotated and labeled large-scale image datasets, ensuring high-quality ground truth for robust model training. Collaborated with engineering teams to improve detection accuracy and deployment readiness.
\end{itemize}

% ---------- Projects ----------
\section*{Projects}

\textbf{Confluex — Automated Literature Review Platform} \hfill Apr 2026 -- Present
\begin{itemize}[leftmargin=1.5em]
    \item Position: AI Engineer(VinGroup Applied AI Talent Program — A20, Team 143).
    \item Built a full-stack academic research workspace that turns a vague research question into a ranked paper library, grounded summaries, and citation-formatted drafts. Designed a multi-agent LangGraph pipeline (Searcher $\rightarrow$ Reader $\rightarrow$ Writer $\rightarrow$ QA) that expands queries, searches Semantic Scholar and arXiv, deduplicates, and ranks results by embedding-based semantic relevance.
    \item Implemented grounded paper/project chat and a Deep Search agent that synthesizes multi-step reports from academic and web (Tavily) sources, with every claim traceable to a source and QA validation flags to prevent hallucinated citations. Built a Writer agent producing IEEE / APA / Chicago drafts with BibTeX export and a section editor.
    \item Live demo: \href{https://a20-app-143.vercel.app}{a20-app-143.vercel.app}
\end{itemize}

\textbf{A Comparative Study of Vietnamese Scene Text Recognition Models} \hfill 2025 -- Jun 2026
\begin{itemize}[leftmargin=1.5em]
    \item Position: Bachelor Thesis Author, University of Science and Technology of Hanoi.
    \item Built and compared two end-to-end Vietnamese scene-text recognition pipelines on the VinText dataset (2{,}000 images, 56K+ annotated instances): a DBNet (ResNet-18) + CRNN (CTC) baseline and a fine-tuned PaddleOCR PP-OCRv3 system (MobileNetV3 detection, SVTR-LCNet recognition). Implemented NFD Unicode normalization for Vietnamese diacritics and a unified evaluation protocol (IoU $\geq$ 0.5 matching, detection/E2E H-mean, CER, WER).
    \item PP-OCRv3 achieved 51.94\% end-to-end H-mean, 8.26\% CER, and 24.03\% WER on the test set, outperforming DBNet+CRNN (32.47\% E2E H-mean) on recognition while both systems reached comparable detection quality ($\sim$63--65\% H-mean). Benchmarked results against published Dictionary-guided Scene Text Recognition baselines and documented reproducible training pipelines in Jupyter notebooks on Kaggle GPU.
    \item Link: \href{https://github.com/TranGiaKhanh283/A-Comparative-Study-of-Vietnamese-Scene-Text-Recognition-Models}{github.com/TranGiaKhanh283/A-Comparative-Study-of-Vietnamese-Scene-Text-Recognition-Models}
\end{itemize}

\textbf{End-to-End Neural Speaker Diarization with Contrastive Learning} \hfill Nov 2025
\begin{itemize}[leftmargin=1.5em]
    \item Position: Data Scientist.
    \item Designed a progressive speaker diarization system from classical clustering baselines to end-to-end neural diarization (EEND). Built an embedding-based pipeline using ECAPA-TDNN embeddings and spectral clustering, achieving 82.1\% frame-level accuracy on VoxConverse. Developed a Conformer-based EEND model with permutation-invariant training (PIT). Introduced contrastive learning pretraining on CallHome data to improve speaker representations. Implemented attention-based speaker pooling and overlap-aware loss weighting.
    \item Achieved 17.54\% Diarization Error Rate (DER) with under 1 hour training time and 0.18× real-time inference on Kaggle GPU. Conducted component and error analysis, showing attention mechanisms outperform memory-based architectures under limited data regimes.
    \item Link:https://github.com/TranGiaKhanh283/fundadata/tree/main
\end{itemize}



\Needspace{6\baselineskip} 
\textbf{Pickleball videos Action Recognition} \hfill Oct 2025 -- Present
\begin{itemize}[leftmargin=1.5em]
    \item Position: Data Scientist, AI Engineer.
    \item Designed a pose-based action recognition pipeline using YOLOv11-Pose + LSTM, focusing on motion dynamics rather than raw pixels. Built a custom dataset of 2,680 video clips covering 4 pickleball techniques with balanced class distributions and data generalization. Developed a keyframe extraction module using pretrained YOLOv11-Pose to normalize variable-length videos into fixed-length sequences as train and test data. Built, trained and finetuned LSTM action classifiers models on different parameters configurations.
    \item Achieving 74\% accuracy on binary classification and 41\% accuracy on 4-class recognition as the best results. Analyzing and concluding for the future improvement: the model requires more diverse and professional data collection, expanding camera angles and techniques, and exploring advanced architectures such as attention and Transformer-based models to improve robustness and accuracy.
    \item Link:https://github.com/TranGiaKhanh283/Projects
\end{itemize}

% ---------- Skills ----------
\section*{Skills}
\textbf{Programming:} Python, SQL, R \\
\textbf{Frameworks \& Tools:} PyTorch, TensorFlow, Microsoft Excel \\
\textbf{Libraries:} Pandas, NumPy, Matplotlib, scikit-learn, Ultralytics \\
\textbf{Platforms:} Google Colab, Kaggle, GitHub, Roboflow, Hugging Face \\
\textbf{Languages:} English (IELTS 7.5), French (TCF Tout Public B1)

% ---------- Certificates ----------
\section*{Certificates}
Google Data Analytics Certificate \hfill 2025 \\
Google Advanced Data Analytics Certificate \hfill 2026

\end{document}
